In this work, an existing neural network architecture was modified for use in data independent physics-based advective modelling of spatiotemporal data and it was used as a state operator in a Kalman filter model.
The obtained neural network trained with a small dataset of satellite measurements of clouds was shown to generalise well for data unrelated to the training data while exhibiting comparable perormance to classic optical flow methods.
With using the Kalman filter with simple cloud data, the neural network was able to model the state evolution and thus fill missing observations.
In these experiments the concept of using neural networks in physical modelling was validated along with using it as a part of a more complex system.
While the data used in numerical experiments was limited to the cloud optical thickness observations, similar modelling could have applications across multiple domains.
The advective modelling was chosen for the generality and easy understanding but similar neural network frameworks could easily be applied also in modelling other phenomena where existing methods are not as well established.

While considerable work remains to follow up with what has been done for this thesis both with development of the neural network and tuning the Kalman filter is needed for better validation of the model, the results were already encouraging.
As the neural network was mostly a reused image recognition architecture with limited adjustments, a valid case could be made that there are considerable low-hanging fruits to improve performance.
For this work, the physics-based interpretability of the neural network model was a primary goal so engineering the network further with performance in mind and with more sophisticated training with extended and better appropriate data could easily provide superior results.
In application specific tasks, neural networks predicting outputs directly without physics-based modelling have gained popularity but lack in interpretability and in problems where mathematical models exist, the methods studied in this work could provide value in constraining the predictions.

In many ways the application of Kalman filtering in this work was limited and could not show its full potential.
While the examples studied here demonstrated its ability fill missing values in partial observations, more sophisticated treatment of nonstaionary inverse problems related to the observation operator in the state-space model could provide value in many complex problems.
For instance, data fusion of multiple data sources along with the methods studied here could contribute to meaningful improvements in accuracy of many problems.

Even with the wide use and theoretical validity of the ensemble based Kalman filters, major drawbacks remain to be handled.
The accurate representation of the uncertainty of state estimates requires handcrafted engineering of the ensembles with multiple unknown parameters.
But as larger state and observation spaces render the classical Kalman filter unfeasible, the ensemble methods offer a valid workaround.
Another alternative for applications similar to the one presented here could be appropriate dimension reduction techniques in order to continue using the classical Kalman filter for higher dimensional problems as well.
In addition, the filtering problems could be reformulated to better represent the actual task. 
As the neural network state operator models only advection at a few control points and big uncertainties are related to the advection vectors, a potential reformulation would be to incorporate the advection parameters to state vectors.

For meteorological applications, prediction problems related to the used cloud optical thickness observations could also be valuable even though less focus was given to that.
With reliable COT forecasts, the ever increasing solar energy production could be estimated that could be valuable information even for short time horizons in the electricity markets given the variable power of renewable energy sources.
For the forecasting problems, the presented model framework could be already used for complete and reliable estimates of the initial state from which with some adjustments probabilstic forecasts could be produced.

In conclusion, the findings of this work could direct further work toward multiple directions.
In the realm of the study of novel methods, similar techniques could be tested in similar modelling problems to the advection model.
On the other hand, this, and similar methods could benefit from futher verification and optimisation to yield better results.
As for the Bayesian filtering problems, the completing and smoothing of historical advective time series or the incorporation of more sophisticated observation models would be a direct continuation of this work.
Finally, the advective prediction problems related to a specific data and applications could be studied with the model of this work providing complete initial states and a predictor model.
